
A webalkalmazás célja a modern tömegközlekedési rendszerek felhasználóbaráttá és hatékonyabbá tétele, különös tekintettel az utazás megtervezésére, jegyvásárlásra, valamint az autóbuszok követésére. Az alkalmazás lehetővé teszi a felhasználók számára, hogy előre megtervezhessék utazásaikat egy adott város vagy régió tömegközlekedési hálózatán belül. A felhasználók megvásárolhatják jegyeiket egy kiválasztott járatra, valamint megtekinthetik a járatok indulási és érkezési időpontjait, városokra és megállókra bontva, biztosítva ezzel a rugalmasságot és a pontos utazási információkat.

Az alkalmazás egy interaktív térképet is kínál, amelyen valós időben nyomon követhetők az autóbuszok jelenlegi pozíciói. Ez a funkció különösen hasznos, amikor a menetrendben váratlan változások történnek, mivel a felhasználók időben értesülhetnek a késésekről vagy útvonal-módosításokról. A valós idejű követés mellett a rendszer biztosítja az egyes buszmegállók pontos földrajzi helyzetét, amelyeket a felhasználók városok vagy útvonalak alapján szűrhetnek, így könnyebben tervezhetik meg, melyik megállóhoz érkezzenek.

A buszokon található POS terminálokon futó különálló alkalmazás gondoskodik a jegyérvényesítésről és az autóbuszok pozíciójának folyamatos megosztásáról a központi szerverrel, így a felhasználók folyamatosan tájékozódhatnak az autóbuszok valós helyzetéről.

Az alkalmazás másik fontos része az adminisztrációs felület, amelyet a busztársaságok használhatnak. Ezen a felületen lehetőség van a menetrendek kezelésére, új útvonalak és megállók felvitelére, valamint a meglévő adatok frissítésére. Továbbá itt kezelhetők az alkalmazottak információi, beleértve a sofőrök és egyéb személyzet adatait, valamint az autóbuszok műszaki és üzemeltetési adatai. Ez az adminisztrációs rendszer biztosítja a hatékony működést és a menetrendek karbantartását, így a busztársaságok könnyedén frissíthetik szolgáltatásaikat és biztosíthatják a pontosságot.

A dolgozat központi témája a webalkalmazás szerver oldali komponenseinek részletes ismertetése, amelyek biztosítják a felhasználói adatok, a jegyvásárlási tranzakciók, a menetrendi információk, valamint a valós idejű autóbusz-követés zökkenőmentes kezelését.

\vspace{1em}
\textbf{Kulcsszavak:} tömegközlekedés, valós idejű követés, webalkalmazás, Rails API